\documentclass[10pt]{article}
\usepackage[T1]{fontenc}
\usepackage{lmodern,microtype,amsmath,amssymb,amsthm,mathtools,booktabs}
\usepackage[margin=0.82in]{geometry}
\usepackage[colorlinks,linkcolor=blue!50!black,urlcolor=blue!50!black,citecolor=blue!50!black]{hyperref}
\usepackage{xcolor}
\newtheorem{theorem}{Theorem}
\newtheorem{lemma}[theorem]{Lemma}
\newtheorem{corollary}[theorem]{Corollary}
\newcommand{\Z}{\mathbb Z}
\newcommand{\F}{\mathcal F}
\newcommand{\Hh}{\mathcal H}
\newcommand{\Nn}{\mathcal N}
\newcommand{\s}{\sigma}
\newcommand{\B}{\mathcal B}
\newcommand{\norm}[1]{\lVert #1\rVert_1}
\newcommand{\coef}[2]{[#1^{#2}]}
\title{Diameter three in every root lattice and the largest anticodes outside ball sections}
\author{Research note}
\date{November 27, 2026}
\begin{document}
\maketitle
\begin{abstract}
We prove the diameter-three case of Kova\v cevi\'c's isodiametric conjecture in every root lattice $A_n$, including classification of the maximizers. Keeping a diametral pair fixed and sampling coordinates reduces the infinite sequence of ranks to four weighted clique bounds on fifteen axes, which we verify in exact arithmetic. In a second theorem we determine, at each fixed diameter $D\ge2$ and all sufficiently large ranks, the largest anticodes that escape every integer-centered ball section. For $D\ge3$ the maximizer consists of a balanced section truncated by one exterior point, together with that point. At diameter two a tetrahedral configuration ties.
\end{abstract}

\section{The extremal problem}

Put $N=n+1$, $\s(z)=\sum_{i=1}^N z_i$ and $A_n=\{z\in\Z^N:\s(z)=0\}$. A diameter-$D$ anticode is a family $\F\subset A_n$ with $\norm{x-y}\le2D$ for every $x,y\in\F$. For any $c\in\Z^N$ write
\[
  \B_j(c)=\{y\in A_n:\norm{y-c}\le j\}.
\]
The center need not belong to $A_n$. Let $\tau_n(D)$ denote the largest size of a diameter-$D$ anticode contained in no $\B_D(c)$, with value zero if no such anticode exists.

Kova\v cevi\'c's conjecture on largest anticodes predicts that a balanced section maximizes cardinality over all anticodes \cite{kov}. The repository gives exact results at small parameters and a volume asymptotic for fixed rank \cite{repo}. Schubr proves the unrestricted conjecture when $n\ge256D^2$, and in every fixed rank for sufficiently large $D$. His Theorem~3.3 bounds $\tau_n(D)$ between constant multiples of $D^2 M_n(D)/n$, where $M_n(D)$ is the size of a balanced section \cite{schubr}. The two results below are unconditional. They extend the lattice problem singled out by the repository; neither uses the zero-location input.

\begin{theorem}[Kova\v cevi\'c's conjecture at diameter three]\label{diam3}
For every $n\ge1$, the largest size of a diameter-three anticode in $A_n$ is
\[
                      M(N)=N+N\binom N2,\qquad N=n+1.
\]
Every maximizer is a section $\B_3(c)$ with $\s(c)=1$ or $-1$.
\end{theorem}

The proposed section has this size. For example, relative to a center of sum $-1$, a displacement of coordinate sum one and norm at most three is either $e_i$ or a positive mass of two and a negative mass of one. In the latter case choose the negative axis and a two-element multiset on the other $N-1$ axes. For $N=8$ the answer is $232$, and for $N=44$ it is $41{,}668$. Schubr's Theorem~3.6 gives the diameter-three result for $n=4,5,6$ and $n\ge44$, complementing the earlier low-dimensional results. Theorem~\ref{diam3} fills ranks $7\le n\le43$, including all equality cases.

\begin{corollary}[Uniform multisets]
A family of $K$-element multisets on $N\ge2$ symbols, any two of which share at least $K-3$ elements counted with multiplicity, has size at most $N+N\binom N2$. Equality can be attained if and only if $K\ge N+1$.
\end{corollary}
\begin{proof}
Multiplicity vectors have sum $K$ and half their $\ell_1$ distance equals $K$ minus the intersection size. Translating this layer to $A_{N-1}$ permits Theorem~\ref{diam3}. In the orientation with positive and negative budgets $(2,1)$, the center has sum $K-1$ and the minimum in coordinate $i$ is $c_i-1$, attained because there is another axis available. Thus this section fits the nonnegative orthant precisely when every $c_i\ge1$, possible exactly when $K\ge N+1$. The reverse orientation gives minima $c_i-2$ and center sum $K+1$; it fits only when $K\ge2N-1$. The equality classification exhausts the possibilities.
\end{proof}

\subsection*{A fixed-rank calculation that controls unbounded rank}

For the upper bound it suffices to consider families with actual diameter three. Indeed if the diameter is $d<3$, choose a pair $x,y$ attaining it and a root $u=e_i-e_j$ which increases $\norm{y-x}$ by two when added to $y-x$. Choose $i$ at a positive and $j$ at a negative coordinate, or any two distinct axes if $x=y$. The new point $y+u$ has distance $d+1$ from $x$ in $A_n$, and distance at most $d+1$ from every previous member. Iterating strictly enlarges the family to actual diameter three.

Translate a diametral pair to $0,h$. The positive and negative parts of $h$ both have mass three, hence each is of type $3$, $21$ or $111$. Coordinate permutations and reflection give six unordered pairs of types. Write $s$ for the support size of $h$. Every member $y$ has $\norm y,\norm{y-h}\le6$. Consequently
\[
 2\sum_{i\notin\operatorname{supp}h}|y_i|
 \le \norm y+\norm{y-h}-\norm h\le6 .
\]
If $k(y)$ counts the distinct nonzero coordinates of $y$ off this support, then $k(y)\le3$. This is what makes a finite verification possible.

Fix $m=15$. For $N\ge m$ keep all $s$ anchor coordinates and choose $m-s$ of the remaining $N-s$ coordinates uniformly. Retain the members $y$ supported on the resulting set of $m$ axes, deleting the other, zero, coordinates. A member with $k$ free axes survives with probability $\binom{m-s}k/\binom{N-s}k$. Thus the expected sum of the weights
\[
 w_k=\frac{\binom{N-s}k}{\binom{m-s}k}
     =\sum_{j=0}^3\binom{N-m}j\,
       \frac{\binom{m-s}{k-j}}{\binom{m-s}k}
 \tag{A}\label{weight}
\]
over the surviving family equals its original cardinality. The expansion is Vandermonde's identity, with a binomial coefficient zero if the lower index is negative. For $j=0,1,2,3$ put
\[
 b_j(k)=\frac{\binom{m-s}{k-j}}{\binom{m-s}k},\qquad
 d_j=\sum_{i=0}^{j}(-1)^{j-i}\binom ji M(m+i).
\]
The four numbers $d_j$ are $1590,346,47,3$, respectively. As $M$ is cubic,
$M(N)=\sum_{j=0}^3\binom{N-m}{j}d_j$.

\begin{lemma}[Finite verification]\label{finite}
For every diameter-three anticode $E\subset A_{m-1}$ containing $0,h$, and every $0\le j\le3$,
\[
                        \sum_{y\in E} b_j(k(y))\le d_j.
\]
For each $2\le N\le15$, every such anchored anticode has size at most $M(N)$, with equality only for a balanced section.
\end{lemma}

The certificate for this lemma is given below. Granting it, apply the four inequalities to every restriction in \eqref{weight}. All multipliers $\binom{N-m}{j}$ are nonnegative, so taking expectations proves $|\F|\le M(N)$ for every $N\ge m$.

If equality holds, the $j=0$ inequality, which has coefficient one, is an equality for every restriction: all deficits are nonnegative and the sampling space is finite. Each restriction is consequently a full balanced section. Any covering center for the two anchors must be at integer $\ell_1$ distance exactly three from each and on their common shortest path, since $\norm h=6$. In particular it vanishes off the fixed $s$ axes. Two sampled coordinate sets differing in one axis share at least three free axes. Suppose their centers were $c\ne c'$. On those free axes take a displacement $v$ of length three with three unit entries whose sum is $-\s(c)$, possible for both balanced signs. The point $c+v$ belongs to the first section and is supported on the intersection, so it belongs to the second restriction as well. But $\norm{c+v-c'}=3+\norm{c-c'}>3$. This contradiction makes the centers equal. Exchanges of single axes connect all the sampled sets. Every original member survives some restriction because $k(y)\le3$, so the original family is contained in a single balanced section and fills it by cardinality. This proves Theorem~\ref{diam3}.

\subsection*{Verification of the finite inequalities}

The file \texttt{check\_diameter\_three.py} is a complete deterministic checker using only the Python standard library. Its output is recorded in \texttt{diameter-three-verification.txt}. For each anchor type it enumerates all possible $y\in A_{N-1}$ other than the anchors with $\norm y,\norm{y-h}\le6$. It processes the axes in order, tracking separately the remaining positive and negative budgets of $y$ and $y-h$. Each of the four budgets starts at three. Extending by every integer within the four budgets and checking coordinate sum zero at the last axis enumerates the lens exactly.

Make these points vertices, adjacent when their $\ell_1$ distance is at most six. A clique, together with the anchors, is an anticode. Signed coordinate levels encode each point as a token set $Y$, so adjacency is equivalent to $|Y\cap Z|\ge(|Y|+|Z|-6)/2$. The checker implements this comparison with integer bitsets. For each $j$ it clears the denominators of the weights $b_j$ by their least common multiple, omits zero-weight vertices and tests for a clique of weight at least one more than the resulting integral bound. The anchors account for two units for $j=0$ and zero for the other indices. The smaller ranks require only the unit-weight test.

For completeness, the search uses the following elementary bound. At a node let $T$ be the weight still needed from a candidate pool. Take an independent set in the pool, and reduce the residual demand of each of its members by the same positive integer $t$, never below zero. This pays at most $t$ units toward any clique, since a clique uses at most one vertex of the independent set. Repeat with a total budget of $T-1$. Every successful clique must meet the set $U$ of vertices with positive residual demand. Branch on the first vertex $v\in U$ it uses, discarding earlier members of $U$ and keeping all remaining neighbors of $v$. The recursive target is $T$ less the \emph{original} weight of $v$. This exhausts the possibilities; if $U$ is empty the node is closed.

For equality the unit-weight search targets $M(N)-2$. It lists every candidate balanced center through the two anchors by enumerating the integer points coordinatewise between $0,h$ with norm three and coordinate sum $\pm1$. Throughout a branch it keeps the centers still compatible with every chosen vertex. If a compatible center covers the remaining pool, the branch is finished. A leaf of target weight with no compatible center would refute the equality assertion. The checker finds none.

Here are the search-node counts at the calibration rank. Counts include the root and depend on the ordering of vertices. All entries represent completed searches.
\begin{center}
\begin{tabular}{lrrrrr}\toprule
Anchor partitions & Size & Equality & $j=1$ & $j=2$ & $j=3$\\\midrule
$3,3$     &7&65&9&37&57\\
$3,21$    &14&56&4&138&163\\
$3,111$   &39&194&22&229&218\\
$21,21$   &35&198&46&284&330\\
$21,111$  &133&329&128&645&505\\
$111,111$ &467&548&555&1858&2411\\\bottomrule
\end{tabular}
\end{center}
The archive also includes an independent C++17 implementation, \texttt{verify\_diameter\_three.cpp}, and its transcript, \texttt{independent\_verification.txt}. It independently rechecks all 90 size or lifting bounds and all 72 equality tests, for 162 successful checks in total.

This finite reduction uses the nonnegative coefficients of the binomial expansion in \eqref{weight}. An ordinary maximum-clique calculation at rank fourteen alone would give only the constant coefficient. Three additional weighted inequalities control the effect of introducing arbitrarily many new axes. For diameter $D$ the same sampling identity has just $D+1$ terms, since the two diameter anchors leave at most $D$ units of mass on free axes. Whether this approach can settle diameters beyond three at a feasible calibration rank remains open.

\section{Anticodes outside every ball section}

Write $P=\lceil D/2\rceil$ and $Q=\lfloor D/2\rfloor$. Choose $c$ with $\s(c)=Q-P$. Choose a vector $h$ with $P+1$ entries $+1$, $Q+1$ entries $-1$ and all remaining entries zero, and put $x=c+h$. Define
\[
 \Hh(c,x)=\{x\}\cup\bigl(\B_D(c)\cap\B_{2D}(x)\bigr).
 \tag{1}\label{model}
\]
Changing the choice of $c$ or of the signed coordinates of $h$ gives an equivalent family, by a translation in $A_n$ and a permutation of coordinates.

\begin{theorem}\label{main}
Fix $D\ge2$. For all sufficiently large $n$, $\tau_n(D)=|\Hh(c,x)|$. If $D\ge3$, every maximizer has the form \eqref{model} up to a permutation, a translation in $A_n$ and global negation. At $D=2$ the value is $4N-6$, and the other maximizing family, up to these symmetries, is
\[
  \{b+e_i:1\le i\le N\}\ \cup\
  \{b+e_i+e_j-e_k:i,j\in\{1,2,3\},\ i<j,\ k\notin\{i,j\}\},
  \qquad \s(b)=-1 .
  \tag{2}\label{tet}
\]
For every fixed $D$ the leading term is
\[
 \tau_n(D)=\frac{P(P+1)+Q(Q+1)}{P!Q!}\,N^{D-1}+O_D(N^{D-2}).
 \tag{3}\label{asympt}
\]
\end{theorem}

For comparison, $M_n(D)=N^D/(P!Q!)+O_D(N^{D-1})$ by counting the points with $P$ distinct positive and $Q$ distinct negative coordinates relative to the center. Thus the leading constant in $N\tau_n(D)/M_n(D)$ is $P(P+1)+Q(Q+1)$.

\section{Extracting centers}

We work with fixed $D$ throughout. The constants in the error terms below depend only on $D$, uniformly over all anticodes and all their centers. Write $r=D-1$.

\begin{lemma}\label{core}
Suppose $\F$ is a diameter-$D$ anticode contained in no $\B_D(c)$. There is a set $K\subset\Z^N$ of $O_D(1)$ centers with the following properties.
\begin{enumerate}
\item All but $O_D(N^{D-2})$ members of $\F$ lie in $\bigcup_{b\in K}\B_r(b)$.
\item Every $y\in\F$ satisfies $\norm{y-b}\le D+1$ for each $b\in K$.
\item Every $b\in K$ satisfies $\s(b)\equiv r\pmod2$, and distinct centers in $K$ have distance exactly two.
\end{enumerate}
\end{lemma}

\begin{proof}
Translate a member of $\F$ to zero. Each $y\in\F$ then has $\norm y\le2D$. Encode it as a set of $\norm y$ tokens by recording sign and level $1,\ldots,|y_i|$ on each nonzero coordinate $i$. Symmetric difference realizes $\ell_1$ distance exactly. The token universe has size at most $4DN$. Separate the family according to token set size $k\le2D$.

Fix $R=6D+2$. Start a tree for each $k$, from the empty token set $C$. At a node consider the encoded members containing $C$ and their residual sets $Y\setminus C$, of common size $j=k-|C|$. If $j\le D-2$, end this branch. It accounts for $O_D(N^{D-2})$ members. If there are $R$ pairwise disjoint residual sets, end the branch and record $C$ and $j$. Otherwise take a maximal disjoint collection. Its union has at most $2DR$ tokens and meets every residual set. Branch by adjoining each of these tokens in turn to $C$. Every member at the node proceeds to some child. The depth and the number of children at each node are bounded solely in terms of $D$.

At a recorded node $C$ is the intersection of the $R$ witnessing members, so it is a valid signed prefix encoding of a vector $b$. For each $y\in\F$ some residual set misses its encoding. Distance to that witness gives
\[
 \norm{y-b}+j\le 2D.
\]
Two witnesses are $2j$ apart, whence $j\le D$. If $j=D$ the preceding inequality places $\F$ in $\B_D(b)$, contrary to hypothesis. All recorded nodes therefore have $j=r$, and all members at such a node are at distance $r$ from $b$. The parity of $\s(b)$ follows from the parity of its distance to any member.

For two recorded nodes select a petal at the first disjoint from the second core, then one at the second disjoint from the first core and the selected petal. Their respective forbidden sets contain at most $4D$ tokens, so the choices exist. The distance between the selected witnesses is $2r+\norm{b-b'}\le2D$. Parity makes the distance of distinct centers at least two. The small terminal nodes together yield the claimed error. Undoing the initial translation preserves all the conclusions.
\end{proof}

\begin{lemma}\label{shape}
A finite set of mutually distance-two points in $\Z^N$ either lies among the unit neighbors of one integer vector, or consists of the four points
\[
 b,\quad b+u+v,\quad b+u+w,\quad b+v+w
 \tag{4}\label{tetcore}
\]
for three signed unit vectors on distinct axes.
\end{lemma}
\begin{proof}
The cases of at most one point are immediate. Translate one point to zero. Every remaining point is a multiset of two signed unit steps, with repeated steps allowed and opposite steps on one axis excluded. Two such multisets give distance two exactly when their intersection has size one. A common step supplies a common unit neighbor. If there is no common step, a double step is impossible. Take two pairs $\{u,v\},\{u,w\}$. A pair missing $u$ must be $\{v,w\}$, and every pair meeting all three is one of these. Opposite steps cannot appear because each pair appears. This gives \eqref{tetcore}.
\end{proof}

For a center $b$ with the parity of $r$, put $u=(r-\s(b))/2$, $v=(r+\s(b))/2$. If both are nonnegative then
\[
  |\B_r(b)|=\frac{N^r}{u!v!}+O_D(N^{r-1}).
  \tag{5}\label{slice}
\]
The leading term counts $u$ distinct positive and $v$ distinct negative steps. All other displacements use at most $r-1$ distinct coordinates. The section is empty if either budget is negative. In particular the largest possible leading coefficient is
$B_r=(\lfloor r/2\rfloor!\lceil r/2\rceil!)^{-1}$.

\section{The construction and its leading coefficient}

The family in \eqref{model} has diameter at most $D$. It contains $\B_r(c+s)$ for every signed coordinate step $s$ of $h$. Indeed the distances of its center $c+s$ to $c$ and $x$ are $1$ and $D+1$. These $D+2$ smaller sections have overlaps of order $O_D(N^{r-1})$. To see this, for distinct centers $b,b'$ in this list, a boundary vector $y-b$ avoiding the at most two coordinates of $b'-b$ has $\norm{y-b'}=r+2$. Boundary vectors touching these specified coordinates number $O_D(N^{r-1})$, as do all interior vectors.

For a positive step of $h$, the budgets for $\B_r(c+s)$ in \eqref{slice} are $(P-1,Q)$; for a negative step they are $(P,Q-1)$. Their combined leading coefficient is
\[
 L_D=\frac{P+1}{(P-1)!Q!}+\frac{Q+1}{P!(Q-1)!}
     =\frac{P(P+1)+Q(Q+1)}{P!Q!}.
 \tag{6}\label{ld}
\]
The model thus has at least $L_D N^r-O_D(N^{r-1})$ points.

For large $N$ it escapes every radius-$D$ section. For each of the $D+2$ centers $b=c+s$, take $3D+2$ disjointly supported displacements of length $r$ with the indicated budgets. If $z$ were a covering center, then $\norm{z-b}\le D+r$ by any one witness. Some chosen displacement avoids the support of $z-b$, whence $\norm{z-b}+r\le D$. Thus $z$ is within one of every $c+s$. Three different signed directions on distinct axes have $c$ as their unique common unit neighbor, so $z=c$. But $\norm{x-c}=D+2$.

\section{Classification of the optimizers}

Apply Lemmas~\ref{core} and \ref{shape} to a maximizer. If the centers form a star, let $c$ be a common unit neighbor. There is an outsider $x\in\F\setminus\B_D(c)$. We may assume $K$ is nonempty since otherwise the family has order $O_D(N^{r-1})$. Since $\s(c)\equiv D\pmod2$ and $\norm{x-b}\le D+1$ for each $b\in K$, one has $\norm{x-c}=D+2$. Every direction $b-c$ must be a signed coordinate step toward $x$, because taking that step reduces the distance to $x$ by one. In particular the directions occur on distinct axes.

Let $A$ and $B$ be the positive and negative masses of $x-c$, so $A+B=D+2$ and $\s(c)=B-A$. There are at most $A$ positive and $B$ negative directions in $K$. Each sign with zero available slice contributes zero. If $A=0$ or $B=0$ all slices in the possible directions are empty. Otherwise set $a=A-1$, $b=B-1$, giving $a+b=D$. The total available coefficient is bounded above by
\[
 \frac{A}{(A-2)!(B-1)!}+\frac{B}{(A-1)!(B-2)!}
 =\frac{a(a+1)+b(b+1)}{a!b!},
 \tag{7}\label{starweight}
\]
where inverse factorials of negative integers mean zero. This bounds $|\F|/N^r$ up to $O_D(1/N)$ by Lemma~\ref{core} and \eqref{slice}.

The right side is uniquely maximized by $\{a,b\}=\{P,Q\}$. Indeed if $a\ge b\ge1$ and $S=a(a+1)+b(b+1)$, replacing $(a,b)$ by $(a+1,b-1)$ multiplies it by $b(S+2(a-b+1))/((a+1)S)<1$, since $S>2b$. At balance every positive and negative direction has positive weight. To achieve $L_D$ all $D+2$ directions must therefore occur, and all entries of $x-c$ must be $0,1,-1$. Any other outsider must be at distance $D+2$ and point in every one of the same $D+2$ directions, so it equals $x$. The family is then contained in \eqref{model}, up to reflection.

In the tetrahedral case the upper coefficient is at most $4B_r$. From \eqref{ld}, $L_D=2(\lfloor D/2\rfloor+1)B_r$, which is larger for $D\ge4$. For $D=3$ equality would require every center of \eqref{tetcore} to have sum zero, by \eqref{slice}. The four sums take two different values since the three signed steps include two of equal sign. Some center thus has coefficient strictly smaller than $B_r$.

For fixed $D$ the above coefficient bounds assume only finitely many values. In the star case they depend on $A,B$ and the two numbers of available directions. In the tetrahedral case only four integral slice budgets occur, with a budget outside $[0,r]$ contributing zero. All errors are uniform by the bounded tree in Lemma~\ref{core}. Consequently any strictly smaller coefficient is excluded for every sufficiently large $N$, by comparison with the model. For $D\ge3$ this leaves only containment in \eqref{model}. Maximal cardinality forces equality, and \eqref{asympt} follows. The finite-coefficient argument also proves there is a positive constant $\delta_D$ such that any family escaping all sections and all copies of \eqref{model} has size at most $(L_D-\delta_D)N^r$ for sufficiently large $N$, when $D\ge3$.

\paragraph{Diameter two.}
Here every tetrahedral center must have sum $\pm1$ to tie $L_2=4$. The four sums in \eqref{tetcore} consist of three equal values and one other, because each signed step has sum $\pm1$. By reflecting and choosing the exceptional center as the base $b$, the sums are $-1,+1,+1,+1$. The other centers are $b+e_i+e_j$ for the three pairs from three axes. Indeed their displacements from $b$ are positive two-step pairs; Lemma~\ref{shape}'s intersection argument gives the three distinct axes. For $y\in\F$, $z=y-b$ has sum one and norm at most three. Norm one gives $z=e_k$. At norm three the positive mass is two, and the support of the positive part must meet each of the three pairs by the constraint $\norm{z-e_i-e_j}\le3$. Thus $z=e_i+e_j-e_k$ with $k\notin\{i,j\}$. This proves containment in \eqref{tet}.

Conversely \eqref{tet} is the union of the four radius-one sections around these centers. The centers are at mutual distance two, so this union has diameter at most two in $A_n$. For large $N$ the disjoint-displacement argument above forces any covering center to be within one of all four centers. They have no such neighbor. The displayed points are distinct and give $N+3(N-2)=4N-6$. In the star model at $D=2$, aside from $c$ and $x$ one may take precisely $c+e_i-e_j$ with $i\ne j$ for which $i$ meets the two positive positions of $h$ or $j$ meets the two negative ones. There are $4(N-1)-4$ such points. The star model lies in a single section after deleting one point. At large $N$ the four sets around the tetrahedral centers still force a covering center within one of each even after one deletion, so the two models are inequivalent. Thus both models have the same size, and the coefficient gap argument yields the asserted classification also at $D=2$. This completes the proof of Theorem~\ref{main}.

\section{Exact evaluation and scope}

Put $G(u,v)=1+u/(1-u)+v/(1-v)$. Its coefficients encode a signed integer coordinate with positive and negative masses marked separately. Every $y$ at distance at most $D-2$ from $c$ is compatible with $x$. A boundary point $y=c+z$ is compatible exactly when some coordinate of $z$ has the sign prescribed by $h$ at a nonzero position of $h$. The exact answer for sufficiently large $n$ is therefore
\[
\begin{split}
 M_n(D)&=\sum_{j=0}^{Q}[u^{P-j}v^{Q-j}]\,G(u,v)^N,\\
 \tau_n(D)&=1+M_n(D)-
 [u^P v^Q]\frac{G(u,v)^{N-D-2}}{(1-v)^{P+1}(1-u)^{Q+1}} .
\end{split}
\tag{8}\label{count}
\]
For example, at $D=3$, $P=2,Q=1$. The balanced section contains $N$ points of norm one relative to $c$, and $N\binom{N}{2}$ of norm three. At norm three fix the negative coordinate. If it is one of the two negative positions of $h$, all positive pairs on the other $N-1$ axes are allowed. At every other negative coordinate keep precisely the positive pairs meeting one of the available positive positions of $h$. This gives
\[
           \tau_n(3)=4N^2-15N+22
           \quad\text{for sufficiently large }N.
\tag{9}\label{dthree}
\]
The accompanying checker independently computes \eqref{count}, enumerates both models at small parameters and checks \eqref{dthree}. The finite enumerations evaluate the proposed models; the proof above establishes optimality in the eventual range.

\begin{corollary}[All multiset sizes at defect three]\label{multi}
For $N\ge48$ and every $K\ge3$, the maximum size of a family of $K$-multisets on $N$ symbols with pairwise intersections of size at least $K-3$ is
\[
 \max\left\{\binom{N+2}3,\quad N+\min(N,K-1)\binom N2\right\}.
 \tag{B}\label{allmultis}
\]
The first term counts all multisets containing a fixed $(K-3)$-multiset. The second counts all multisets sharing at least $K-2$ entries with a fixed $(K-1)$-multiset of maximum support. The larger construction is extremal. The first term wins for $3K<N+7$, the second for $3K>N+7$, and their sizes agree at the boundary.
\end{corollary}

\begin{proof}
Work in the nonnegative layer of coordinate sum $K$. The first construction has $S=\binom{N+2}3$ members. For the second take a nonnegative center of sum $K-1$ and support size $\min(N,K-1)$. Positive and negative displacements have budgets $(2,1)$. There are $N$ norm-one displacements; for norm three, choose a negative axis in the support and a positive two-multiset on the other $N-1$ axes. This gives $H=N+\min(N,K-1)\binom N2$.

An optimal family can be chosen in a radius-three section. For sufficiently large $N$ this follows immediately from Theorem~\ref{main}, since $\tau_{N-1}(3)=O(N^2)$ while $S$ grows cubically, uniformly in $K$. Here is the explicit threshold. For $K=3$ every multiset is allowed. For $4\le K\le N-3$ the multiset complete-intersection theorem of F\H{u}redi, Gerbner and Vizer \cite[Theorem 1.8]{fgv}, and the realization in \cite[Theorem 1.4 and the subsequent display]{mp}, give a maximizing family requiring intersection at least $t+j$ with a fixed $(t+2j)$-set, where $t=K-3$ and $0\le j\le3$. In vector coordinates this is a radius-three section since $K+(t+2j)-2(t+j)=3$. For $K\ge N-2$, Schubr's Proposition~3.5 \cite{schubr} bounds the size of any family outside all such sections by
\[
 \Phi(3N,3)=1+\sum_{i=0}^3\binom{3N}{i}-\binom{3N-4}3
           =(45N^2-69N)/2+22 .
\]
The lower bound $H\ge N+(N-3)\binom N2$ exceeds this when $N=48+z$, $z\ge0$: the difference of the latter lower bound and $\Phi$ equals $z^3/2+95z^2/2+1141z+602$.

It remains to optimize sections. Clamping any negative center entries to zero uses up at least one unit of radius for all points in the orthant. The resulting radius-two ball has at most $1+2N+2N^2<S$ points at $N\ge48$. The same bound holds if the parity of the center sum makes the usable radius at most two. For nonnegative centers with the required parity, the full positive and negative budgets are $(3,0),(2,1),(1,2),(0,3)$ whenever the section is nonempty. The first and last cases have at most $S$ members, and the second has at most $H$ by the preceding count.

For $(1,2)$ let $v$ be the number of positive entries of the center and $u$ the number at least two. Its sum is $K+1$. The size is
\[
 v+(N-2)\binom v2+(N-1)u .
\]
These terms correspond to one negative unit, distinct negative axes, and a repeated negative axis. For fixed $v$ the maximum uses $u=\min(v,K+1-v)$. Increasing $v\ge1$ by one when possible gains at least $1+(N-2)v-(N-1)\ge0$, so one may take $v=\min(N,K+1)$. If $K\le N-1$, subtracting the resulting size from $H$ gives $(K-1)(N-2)(N-K-1)/2\ge0$. At $K=N$ the difference is $(N-2)(N-1)/2$. For $K\ge N+1$, $H=M(N)$ is the unrestricted maximum by Theorem~\ref{diam3}. This proves \eqref{allmultis}, and comparing its terms gives the transition.
\end{proof}

Meagher and Purdy proposed in their final section \cite{mp} that an extremizer can be obtained by requiring intersection at least $t+j$ with one $(t+2j)$-multiset, and specified a preferred support and distribution of multiplicities. Corollary~\ref{multi} gives the proposed extremal value and realizes the construction at $t=K-3$ for every $N\ge48$. The two displayed constructions take $j=0,1$; their center multiplicities can be chosen as evenly as possible over maximum support. The range $K>N-3$ lies beyond the hypothesis $N\ge 2K-t$ of the classical theorem.

Two questions remain. What is the least rank at which the nonsection classification stabilizes for each $D$? Can the same exact families be proved extremal under a uniform hypothesis of the form $n\ge C D^2$? Neither follows from the fixed-$D$ coefficient comparison. The unrestricted isodiametric conjecture in the remaining ranks and diameters is a separate question.

\begin{thebibliography}{9}
\bibitem{fgv} Z. F\H{u}redi, D. Gerbner, and M. Vizer. A discrete isodiametric result: The Erd\H{o}s--Ko--Rado theorem for multisets. \emph{European Journal of Combinatorics} 48 (2015), 224--233. \href{https://doi.org/10.1016/j.ejc.2015.02.023}{doi:10.1016/j.ejc.2015.02.023}.
\bibitem{mp} K. Meagher and A. Purdy. Intersection theorems for multisets. \emph{European Journal of Combinatorics} 52 (2016), 120--135. \href{https://arxiv.org/abs/1504.06657}{arXiv:1504.06657}.
\bibitem{kov} M. Kova\v cevi\'c. \emph{An Isodiametric Theorem and Lattice Diameter-Perfect Codes in $A_3$}. 2026. \href{https://arxiv.org/abs/2607.21037}{arXiv:2607.21037}.
\bibitem{repo} V. Reshetnikov, \emph{ProveIt}, Analysis/GeneralizedRiemannHypothesis, Article, section on coefficient collisions and higher-dimensional anticodes. \href{https://github.com/VladimirReshetnikov/ProveIt/blob/main/Analysis/GeneralizedRiemannHypothesis/Article/sections/gpp.tex}{Repository source}.
\bibitem{schubr} S. Schubr. \emph{Maximum anticodes and balanced-anticode tilings in the root lattices $A_n$}. Luminary Research Archive, edition 1, November 2026, Theorems 3.2--3.4. \href{https://research.luminaryoddities.org/maximum-anticodes-balanced-anticode-tilings-root-lattices/v/1/files/Anticodes_Root_Lattices.pdf}{Author's paper}.
\end{thebibliography}
\end{document}
